% !TEX root = ../matan.tex

\section{Гильбертово пространство}

\begin{Definition}{Гильбертово пространство}{}
	\textbf{Гильбертово пространство} $\mathcal{H}$~--- это линейное пространство над $\RR$, на котором задано скалярное произведение
	\[
	\langle \cdot, \cdot \rangle \colon \mathcal{H} \times \mathcal{H} \to \RR,
	\]
	и которое полно относительно нормы, порождённой этим скалярным произведением:
	\[
	\|x\| := \sqrt{\langle x, x \rangle}.
	\]
	
	Для всех $x,y,z \in \mathcal{H}$ и $\lambda \in \RR$ выполнены свойства:
	\begin{enumerate}
		\item $\langle x, x \rangle \ge 0$, причём $\langle x, x \rangle = 0 \Leftrightarrow x = 0$;
		\item $\langle \lambda x, y \rangle = \lambda \langle x, y \rangle$ и $\langle x, \lambda y \rangle = \lambda \langle x, y \rangle$;
		\item $\langle x + y, z \rangle = \langle x, z \rangle + \langle y, z \rangle$;
		\item $\langle x, y \rangle = \langle y, x \rangle$.
	\end{enumerate}
\end{Definition}

Скалярное произведение автоматически задаёт норму $\|x\| = \sqrt{\langle x, x \rangle}$.

\begin{Proposition}{Неравенство Коши--Буняковского--Шварца}{}
	Пусть $\mathcal{H}$~--- гильбертово пространство. Тогда для любых $x,y \in \mathcal{H}$ выполнено
	\[
	|\langle x, y \rangle| \le \|x\| \cdot \|y\|.
	\]
\end{Proposition}

\begin{proof}
	
	$\lambda \in \RR$
	
	\[
	0 \le \langle x + \lambda y, x + \lambda y \rangle
	= \langle x,x \rangle + 2\lambda \langle x,y \rangle + \lambda^2\|y\|^2.
	\]
	
	Дискриминант квадратного трёхчлена по $\lambda$ не больше нуля, т.е.
	
	\[
	4(\langle x,y \rangle)^2 - 4\langle x,x \rangle\langle y,y \rangle \le 0.
	\]
	
	Следовательно,
	\[
	|\langle x,y \rangle| \le \|x\|\cdot\|y\|.
	\]
\end{proof}

\begin{Proposition}{Неравенство треугольника}{}
	Для любых $x,y \in \mathcal{H}$ выполнено
	\[
	\|x+y\| \le \|x\| + \|y\|.
	\]
\end{Proposition}

\begin{proof}
	По определению нормы и по неравенству Коши--Буняковского--Шварца:
	\begin{align*}
		\|x+y\|^2
		&= \langle x+y, x+y \rangle \\
		&= \|x\|^2 + 2\langle x,y \rangle + \|y\|^2 \\
		&\le \|x\|^2 + 2|\langle x,y \rangle| + \|y\|^2 \\
		&\le \|x\|^2 + 2\|x\|\|y\| + \|y\|^2 \\
		&= (\|x\| + \|y\|)^2.
	\end{align*}
	Так как обе части неотрицательны, получаем требуемое неравенство.
\end{proof}

\begin{Remark}{}{}
	
	Естественнее было бы рассматривать функции
	
	\[ L^2([a,b]) = \left\{f \st [a,b] \rightarrow \RR, \quad \int_a^b |f(x)|^2\,dx < +\infty \right\} \]
	
	\[ \|f\|^2_{L^2} := \int_a^b |f(x)|^2\,dx \]
	
	но $\int_a^b |f(x)|^2\,dx$ может оказаться равным $0$ для $f \neq 0$ как функции, если различать значения в отдельных точках.
	
	Отложим до теории меры
\end{Remark}

\begin{Remark}{}{}
	
	Далее считаем, что все наши гильбертовы пространства~--- сепарабельные.
	
	То есть существует счётное всюду плотное подмножество. (подмножество, замыкание которого равно всему множеству)
	
	Т.е. $\forall \, \mathcal{H} \, \exists \,  E \subset \mathcal{H}$ (счётное) такое что $\forall \,  x \in \mathcal{H}  \, \exists \,  \{ x_k \} \subset E \st  \|x_k - x\| _{\mathcal{H}} \longrightarrow 0 (n \rightarrow +\infty)$
\end{Remark}

\begin{Definition}{Сходимость ряда в гильбертовом пространстве}{}
	$\mathcal{H}$~--- гильбертово пространство, $\{x_k\} \subset \mathcal{H}$
	
	Считаем, что $x_n \longrightarrow x \Leftrightarrow  \|x_n - x\| _{\mathcal{H}} \longrightarrow 0 \quad (n \rightarrow +\infty)$
	
	Ряд $\sum_{k=1}^\infty x_k$ сходится в $\mathcal{H}$ к $S \in \mathcal{H}$, если
	
	\[  \|\sum_{k=1}^Nx_k - S\|  \longrightarrow 0 \quad (N \rightarrow +\infty) \]
\end{Definition}

\begin{Definition}{}{}
	$\mathcal{H}$~--- гильбертово пространство, $x, y \in \mathcal{H} \st x \perp y \Leftrightarrow  \langle x, y \rangle  = 0 $
	
	Система \textbf{ненулевых} векторов $\{x_\alpha\}_{\alpha \in A}$ является \textbf{ортогональной}, если $x_\alpha \perp x_\beta$ для всех $\alpha \neq \beta$.
\end{Definition}

\begin{Proposition}{}{}
	Векторы в ортогональной системе линейно независимы, т.е. если $\sum_{k=1}^N\lambda_k x_k = 0$, где $x_k \in \{x_\alpha\}_{\alpha \in A}$~--- система, то $\lambda_k = 0$ для всех $k$
\end{Proposition}

\begin{proof}
	Зафиксируем $j \in \{1,\ldots,N\}$ и возьмём скалярное произведение с $x_j$:
	\[
	0
	= \left\langle \sum_{k=1}^{N} \lambda_k x_k, x_j \right\rangle
	= \sum_{k=1}^{N} \lambda_k \langle x_k, x_j \rangle
	= \lambda_j \|x_j\|^2.
	\]
	Так как $x_j \neq 0$, то $\|x_j\|^2 > 0$, значит $\lambda_j = 0$. Поскольку $j$ произволен, все коэффициенты равны нулю.
\end{proof}

\begin{Definition}{}{}
	Система $\{x_\alpha\}_{\alpha \in A}$~--- полная, если  $ \operatorname{clos} \; \operatorname{span}\{x_\alpha\}_{\alpha \in A} = \mathcal{H} $ 	т.е. 
	\[ 
	\operatorname{clos} \{ \sum_{k=1}^N\lambda_kx_k \st \lambda_k \in \RR, x_\alpha \in \{x_\alpha\}_{\alpha \in A} \} =  \mathcal{H}
	\]
	
	Полная ортогональная система~--- $\textbf{базис}$ в  $\mathcal{H}$
\end{Definition}

\begin{Example}{}{}
	Базис в $l^2$: $x_k = (0, \ldots, 0, 1, 0, \ldots)$
\end{Example}

\begin{Definition}{}{}
	Базис называется ортонормированным, если $\|x_\alpha\| = 1$ для всех $\alpha$.
\end{Definition}

\begin{Example}{}{}
	Рассмотрим пространство непрерывных функций на отрезке $[0, 1] - C([0, 1])$
	
	$ \langle f, g \rangle  = \int^1_0 f(x)g(x)dx$
	
	$f_0(x) := 0, \qquad f_k(x) := \sin{2\pi kx}, \qquad k \in \NN,  x \in [0, 1]$
	
	$g_0(x) := 1, \qquad g_k(x) := \cos{2\pi kx}, \qquad k \in \NN,  x \in [0, 1]$
	
	\[ f_k \perp f_j; j \neq k \rightarrow \int^1_0 \sin{2\pi kx} \cdot \sin{2\pi jx}dx = 0 \]
	\[ g_k \perp g_j; j \neq k \rightarrow \int^1_0 \cos{2\pi kx} \cdot \cos{2\pi jx}dx = 0 \] 
	\[ f_k \perp g_j; \forall  j, k \rightarrow \int^1_0 \sin{2\pi kx} \cdot \cos{2\pi jx}dx = 0  \quad \forall \, k, j \]
	
	\[ \int_0^1 \sin^2(2\pi kx)\,dx = \frac{1}{2} \]
	
	\[ \operatorname{clos}\,\operatorname{span}\{f_k,g_k\} \supset C([0, 1]) \]
\end{Example}

\begin{Proposition}{}{}
	$\mathcal{H}$ ~--- сепарабельное гильбертово $\Rightarrow \{x_\alpha\}_{\alpha \in A}$~--- ортогональная система, не более чем счётная 
\end{Proposition}

\begin{proof}
	
	WLOG $\{x_\alpha\}_{\alpha \in A}$~--- ортонормированная система
	
	Рассмотрим норму разности:
	
	\[  \|x_\alpha - x_\beta\|  = \sqrt{ \langle x_\alpha, x_\alpha \rangle  - 2 \langle x_\alpha, x_\beta \rangle  +  \langle x_\beta, x_\beta \rangle } = \sqrt{2} \]
	
	Рассмотрим систему шаров $\{ B_\alpha \} \st$
	
	\[ B_\alpha = \{ x \st  \|x - x_\alpha\|   <  \frac{1}{2} \}\]
	
	$B_\alpha \cap B_\beta = \varnothing, \qquad \alpha \neq \beta$
	
	Т.к. $\mathcal{H}$~--- сепарабельная $\Rightarrow  \exists  E = \{ y_k \}_{k \in \NN} \st \operatorname{clos} E =  \mathcal{H}$
	
	Так как $E$ всюду плотно, то $\forall \alpha$ существует $y_{n(\alpha)} \in B_\alpha$. Следовательно,
	
	Таких шаров $B_\alpha$ не более чем счётно
\end{proof}

\begin{Theorem}{}{}
	Пусть $\mathcal{H}$~--- сепарабельное гильбертово пространство
	
	$\{ f_k\}_{k \in \NN}$~--- какая-то линейно независимая система векторов в $ \mathcal{H},  f_k \neq 0$
	
	Тогда $ \exists  \{ \vphi_k\}_{k \in \NN}$
	
	\begin{enumerate}
		\item $\{ \vphi_k\}$~--- ортонормированная система в $\mathcal{H}$
		\item $\vphi_k = \sum_{j=1}^k \lambda_{jk} f_j$~--- линейная комбинация первых $k$ векторов из $\{ f_k\}$
		\item $f_k = \sum _{j=1}^k \tilde{\lambda}_{jk} \vphi_j$~--- линейная комбинация первых $k$ векторов из $\{ \vphi_k \}$
	\end{enumerate}
\end{Theorem}

\begin{proof}
	
	База: 
	$\vphi_1 := \frac{f_1}{\|f_1\|}$
	Переход:
	
	Пусть $\vphi_1, \ldots, \vphi_{k-1}$ удовлетворяют $1-3$
	
	\textbf{Хотим:} получить $\vphi_k$ из $f_1, \ldots, f_{k-1}, f_k$, так, чтобы выполнились $1-3$
	
	3): означает, что 
	\[ f_k = \sum_{j=1}^k\tilde{\lambda}_{jk} \vphi_j = \tilde{\lambda}_{1k}\cdot \vphi_1 + \ldots + \tilde{\lambda}_{(k-1), k}\cdot \vphi_{(k-1)} + g_k \]
	
	Если бы такое представление было, то 
	\[
	\langle f_k, \vphi_j \rangle  = \tilde{\lambda}_{1k} \langle \vphi_1, \vphi_j \rangle  + \ldots + \tilde{\lambda}_{j, k} \langle \vphi_j, \vphi_j \rangle  + \ldots + \tilde{\lambda}_{(k-1), k}\cdot  \langle \vphi_{(k-1)}, \vphi_j \rangle +  \langle g_k, \vphi_j \rangle =
	\]
	
	\[
	= 0 + 0 + \ldots + \tilde{\lambda}_{jk}  \|\vphi_j\|^2 + 0 + \ldots + ( \langle g_k, \vphi_j \rangle  = 0) 
	\]
	
	$ \langle g_k, \vphi_j \rangle  = 0$, т.к. ожидаем, что $\vphi_k \perp \vphi_j$
	
	Итак, если представление есть, то
	
	\[ \tilde{\lambda}_{jk} = \langle f_k, \vphi_j \rangle, \qquad j < k \]
	
	Положим: 
	
	\[ g_k := f_k - \sum_{j=1}^{k-1}\tilde{\lambda}_{j, k}\cdot \vphi_j \]
	
	\[ \vphi_k := \frac{g_k}{ \|g_k\| } \]
	
	При этом $g_k \neq 0$, так как система $\{f_k\}$ линейно независима
	
	$\tilde{\lambda}_{k, k} :=  \|g_k\|  \neq 0$.
	
	Кандидаты на коэффициенты есть. Проверим 1-3.
	
	3. $f_k = \sum_{j=1}^k \tilde{\lambda}_{j,k} \cdot \vphi_j$~--- прямо по определению
	
	1. $\{ \vphi_1, \ldots, \vphi_k \}$ ортогональная система
	
	$\{ \vphi_1, \ldots, \vphi_{k-1} \}$~--- ортонормированная система, $  \|\vphi_k\|  = 1$ по построению
	
	$\langle \vphi_k, \vphi_{j_0} \rangle$, где $j_0 < k$.
	
	\[
	\langle g_k, \vphi_{j_0} \rangle = \left\langle f_k - \sum_{j=1}^{k-1} \tilde{\lambda}_{j,k} \vphi_j, \vphi_{j_0} \right\rangle =
	\]
	\[
	= \langle f_k, \vphi_{j_0} \rangle - \tilde{\lambda}_{j_0,k} \langle \vphi_{j_0}, \vphi_{j_0} \rangle = 0
	\]
	
	2. $\vphi_k = \sum_{j=1}^k \lambda_{j, k} f_j$~--- также ясно, т.к.
	
	\[ \vphi_k = \frac{f_k - \sum_{j=1}^{k-1}\tilde{\lambda}_{j,k} \vphi_j}{ \|f_k - \sum_{j=1}^{k-1}\tilde{\lambda}_{j,k}\cdot \vphi_j\| } =  \text{линейная комбинация из}  f_i ,  1 \le i \le k\]
\end{proof}

\textbf{Далее все Гильбертовы пространства будут полными}

Хочется проверить, что если

$\{ e_k \}_{k \in \NN}$~--- базис в $\mathcal{H}$, то $ \forall x \in \mathcal{H}$~ есть разложение:

\[ x = \sum_{k=1}^\infty c_k e_k, \qquad c_k \in \RR \]

$ \langle x, e_j \rangle  = \sum_{k=1}^\infty c_k  \langle e_k, e_j \rangle  = c_j \cdot \langle e_j, e_j \rangle = c_j$, если $\{ e_k \}$~--- ОНС

\begin{Definition}{}{}
	Пусть $\{ e_k\}$~--- ортонормированная система в $ \mathcal{H}, \qquad x \in \mathcal{H}$
	
	$c_k =  \langle x, e_k \rangle $~--- коэффициент Фурье элемента $x$, по системе $\{ e_k\}_{k \in \NN}$
	
	Ряд (формальный) $\sum_{k=1}^{\infty} c_k e_k$~--- ряд Фурье элемента $x$.
	
	
\end{Definition}

Хочется, чтобы $\sum_{k=1}^\infty c_k e_k = x, \text{если}  \{ e_k \}$~--- ортонормированный базис в $\mathcal{H}$

\begin{proof}
	$N$ элементов~--- $e_1, \ldots, e_N$. Рассмотрим линейную комбинацию:
	
	\[ \sum_{k=1}^N \lambda_k e_k, \qquad \lambda_1, \ldots, \lambda_N \in \RR \]
	
	Минимальное расстояние:
	
	\[  \|x - \sum_{k=1}^N \lambda_k e_k\|  = dist_{\mathcal{H}}(x, \operatorname{span}\{ e_k\}_{k=1}^N) \]
	
	\[  \|x - \sum_{k=1}^N \lambda_k e_k\| ^2 =  \langle x - \sum_{k=1}^N \lambda_k e_k, x - \sum_{j=1}^N \lambda_j e_j \rangle  = \]
	
	\[ = \langle x,x \rangle - 2\left\langle x, \sum_{k=1}^N \lambda_k e_k \right\rangle + \left\langle \sum_{k=1}^N \lambda_k e_k, \sum_{j=1}^N \lambda_j e_j \right\rangle = \]
	
	\[ = \|x\|^2 - 2\sum_{k=1}^N \lambda_k c_k + \sum_{k,j=1}^N \lambda_k\lambda_j \langle e_k,e_j \rangle = \]
	
	\[ =  \|x\|^2 - 2 \sum_{k=1}^N\lambda_k c_k + \sum_{k=1}^N \lambda_k^2 \]
	\[ =  \|x\|^2 - \sum_{k=1}^N c_k^2 + (\sum_{k=1}^Nc_k^2 - 2\sum_{k=1}^N\lambda_k c_k + \sum_{k=1}^N \lambda_k^2) = \] 
	\[ = \|x\|^2 - \sum_{k=1}^N c_k^2 + \sum_{k=1}^N(\lambda_k - c_k)^2. \]
	
	
	
	Минимум: $\lambda_j = c_j, \qquad 1 \le j \le N$, т.е.
	
	\[ dist^2(x, \operatorname{span}\{e_k\}_{k=1}^N) =  \|x\|^2 - \sum_{k=1}^Nc_k^2 \ge 0  \forall  N \]
	\[ \sum_{k=1}^N c_k^2 \le  \|x\|^2 \qquad \forall N  \Rightarrow \sum_{k=1}^\infty c_k^2 \] ~--- сходится и по неравенству Бесселя:
	
	\[ \sum_{k=1}^\infty c_k^2 \le  \|x\|^2 \]
	
	Пусть теперь $\{ e_k\}_{k \in \NN}$~--- ортонормированный базис, т.е. $\operatorname{clos}  \operatorname{span} \{ e_k \} = \mathcal{H} \Leftrightarrow$
	
	\[ \forall  x \in \mathcal{H} \qquad \exists  \{\lambda_k\} \st \sum_{k=1}^N \lambda_k e_k \rightarrow x (N \rightarrow \infty) \] т.е.:
	\[  \|x - \sum_{k=1}^N \lambda_k e_k\| ^2 \longrightarrow 0 \qquad \|x\|^2 - \sum_{k=1}^N c_k^2 \longrightarrow 0 \qquad (N \rightarrow \infty). \]
\end{proof}

\begin{Proposition}{}{}
	Если $\{ e_k\}_{k \in \NN}$~--- ортонормированный базис, то верно равенство Парсеваля:
	
	\[ \sum_{k=1}^\infty c_k^2 = \|x\|^2, \qquad c_k = \langle x, e_k \rangle  \]
\end{Proposition}


Итак: $\mathcal{H}$~--- сепарабельное гильбертово пространство; $\{ e_k \}$~--- ортонормированный базис в $\mathcal{H}$

\[ \mathcal{H} \ni x \longrightarrow \{ c_k\}_{k=1}^\infty, \qquad c_k = \langle x, e_k \rangle\]


\begin{Theorem}{Рисс-Фишер}{}
	Пусть $\{ e_k\}$~--- ортонормированный базис в $\mathcal{H} \st c = (c_1, \ldots, c_k, \ldots) \in l^2 \Rightarrow$
	
	$\exists !  x \in \mathcal{H} \st c_k = \langle x, e_k \rangle_{\mathcal{H}}$ и
	
	$ \|c\| _{l^2} =  \|x\| _{\mathcal{H}}$
\end{Theorem}

\begin{proof}
	
	Конечно же $x = \sum_{k=1}^\infty c_k e_k$
	
	Проверим, что ряд сходится, т.е. последовательность частичных сумм куда-то сходится
	
	Проверим, что это последовательность Коши, т.е. 
	
	\[ \forall \eps > 0 \ \exists N(\eps) \st \forall N \ge N(\eps),\ \forall m \in \NN: \]
	
	\[  \|S_{N+m} - S_N\|  < \eps \]
	
	\[ \|S_{N+m}-S_N\|^2 = \left\|\sum_{k=N+1}^{N+m} c_k e_k\right\|^2 = \sum_{k=N+1}^{N+m} c_k^2. \]
	
	но $c = (c_1, \ldots, c_k, \ldots) \in l^2, \qquad  \|c\|^2 < +\infty = \sum_{k=1}^\infty c_k^2$
	
	\[ \forall  \eps > 0  \exists  N(\eps) \st \forall  m \in \NN \st N \ge N(\eps),  \text{верно}  \sum_{k=N+1}^{N+m}c_k^2 < \eps \]
	
	т.е. ряд $\sum_{k=1}^\infty c_k e_k$ сходится, обозначим эту сумму за $x$
	
	
	Осталось проверить, что $ \langle x, e_k \rangle  = c_k$ и единственность
	
	\[  \langle x, e_k \rangle  \overset{def}{=}  \langle x - S_N, e_k \rangle _{\mathcal{H}} +  \langle S_N, e_k \rangle _{\mathcal{H}} \qquad \forall  N \]
	
	\[ \langle S_N, e_k \rangle = \left\langle \sum_{j=1}^N c_j e_j, e_k \right\rangle = c_k, \qquad N \ge k. \]
	
	\[ | \langle x - S_N, e_k \rangle | \le  \|x - S_N\|\cdot\| e_k \|  \longrightarrow 0 (N \rightarrow +\infty) \]
	
	Единственность:
	
	Пусть $\langle x,e_k \rangle = \langle y,e_k \rangle = c_k$ для всех $k$. Тогда
	
	$(x - y) \perp e_k  \forall  k \Rightarrow (x - y) = 0$ т.к. $\{e_k\}$~--- ортонормированный базис
	
	Пусть $(x-y) = z \neq 0 \Rightarrow c'_k :=  \langle z, e_k \rangle  \neq 0$ (не равняется тождественно)
	
	\[ \operatorname{dist}^2(z, \operatorname{span}\{e_k\}_{k=1}^N) = \|z\|^2 - \sum_{k=1}^N (c'_k)^2. \]
	
	Так как $\{e_k\}$~--- базис, расстояние от $z$ до $\operatorname{span}\{e_k\}_{k=1}^N$ стремится к нулю. Но все коэффициенты $c'_k$ равны нулю, значит это расстояние равно $\|z\|$. Следовательно, $z=0$.
\end{proof}


\begin{Corollary}{}{}
	Любое сепарабельное гильбертово пространство $\mathcal{H}$ изоморфно $l^2$, т.е. есть биекция $\mathcal{A} \st \mathcal{H} \rightarrow l^2$:
	
	\[  \|\mathcal{A}x\|_{l^2} = \|x\|_{\mathcal{H}} \]
	
	$\mathcal{A}$~--- линейный оператор:
	
	$ \langle \mathcal{A}x, \mathcal{A}y \rangle _{l^2} =  \langle x, y \rangle _{\mathcal{H}}$
\end{Corollary}

\begin{proof}
	
	Проверим линейность оператора $\mathcal{A}$:
	
	В $\mathcal{H}$~ есть ортонормированный базис $\{e_k\}_{k=1}^\infty \st \mathcal{A}x = c = \{ c_k \}$
	
	$ \langle \mathcal{A}x, \mathcal{A}y \rangle _{l^2} \overset{\text{def}}{=} \sum_{k=1}^\infty c_k d_k$
	
	\[ 
	S_N = \sum_{k=1}^N c_k e_k \qquad \tilde S_N = \sum_{k=1}^N d_k e_k
	\]
	
	\[
	\begin{aligned}
		\langle x,y \rangle_{\mathcal H}
		&= \langle x-S_N, y-\tilde S_N \rangle
		+ \langle x-S_N, \tilde S_N \rangle \\
		&\quad + \langle S_N, y-\tilde S_N \rangle
		+ \langle S_N, \tilde S_N \rangle.
	\end{aligned}
	\]
	
	Первые три слагаемых стремятся к нулю при $N \to \infty$ по неравенству Коши--Буняковского--Шварца, а
	
	\[  \langle S_N, \tilde S_N \rangle  =  \left\langle \sum_{k=1}^N c_k e_k, \sum_{j=1}^N d_j e_j \right\rangle  = \sum_{k=1}^N c_k d_k. \]
	
	Поэтому $ \langle x,y \rangle_{\mathcal H} = \sum_{k=1}^\infty c_k d_k. $
	
\end{proof}

\begin{Remark}{}{}
	Единичный шар $\overline{B_d(0, 1)} = \{ x \in l^2 :  \|x\| _{l^2} \le 1 \}$~--- не компактен.
\end{Remark}

Recall that:

$X$~--- банахово пространство, если $X$~--- линейное нормированное пространство, полное относительно метрики, порождаемой своей нормой

$\sum_{k=1}^\infty \lambda_k x_k$~--- сходится в $X$, если $ \exists  x \in X \st  \|x - \sum_{k=1}^N\lambda_k x_k\| _X \longrightarrow 0 (N \rightarrow \infty)$

\begin{Definition}{}{}
	Пусть $X, Y$~--- линейные векторные пространства.
	
	Отображение $\mathcal{A} \colon X \rightarrow Y$~--- \textbf{линейный оператор}, если оно линейное:
	
	\[ \mathcal{A}(x + \lambda y) = \mathcal{A}x + \lambda \mathcal{A}y, \qquad \forall  x, y \in X, \lambda \in \RR \]
	
	Если $Y = \RR$, то линейный оператор называется \textbf{линейным функционалом}.
\end{Definition}

\begin{Definition}{}{}
	Пусть $X, Y$~--- банаховы. $\qquad \mathcal{A} \colon X \rightarrow Y$~--- линейный оператор
	
	\[ \mathcal{A} \colon \operatorname{Dom} \mathcal{A} \rightarrow \operatorname{Range} \mathcal{A} \]
	где $\operatorname{Dom} \mathcal{A}$~--- линейное пространство X, а $\operatorname{Range} \mathcal{A}$~--- линейное подпространство $Y$, т.е.
	
	если $x, y \in \operatorname{Dom} \mathcal{A} \Rightarrow x + \lambda y \in \operatorname{Dom} \mathcal{A},  \forall  \lambda \in \RR$. Аналогично для $\operatorname{Range} \mathcal{A}$. 
\end{Definition}

\begin{Example}{}{}
	$X = C([0, 1])$
	
	\[
	\|f\| _X := \sup_{x \in [0,1]} |f(x)|
	\]
	
	\[
	\mathcal{A}f := f', \qquad \operatorname{Dom}\mathcal{A} = \{ f \in C([0,1]) \st f' \text{ существует на } (0,1) \}
	\]
	
\end{Example}

Зададим норму оператора $\mathcal{A}$:

\[ \|\mathcal{A}\|:= \sup_{\|x\|_X \le 1}\|\mathcal{A}x\|_Y \]

Оператор называется ограниченным, если $ \|\mathcal{A}\|   <  +\infty$

\begin{Example}{}{}
	$(x^k)' = kx^{k-1}$, поэтому при $x=1$ получаем $(x^k)'(1)=k$.
	
	\[
	\|x^k\| _{C([0, 1])} = \max_{0 \le x \le 1} |x^k| = 1
	\]
	
	Пример ограниченного оператора:
	
	$(\mathcal{A}f)(x) := f(x)g(x), \qquad \text{где }  g \in C([0,1])$~--- фиксирована
	
	\[  \|\mathcal{A}\|  := \sup_{\|f\|_{C([0,1])} \le 1}  \|\mathcal{A}f\|  = \sup_{x \in [0, 1], f \in C([0,1 ])} |f(x)g(x)| \le \]
	
	\[ \sup |f(x)| \cdot \sup |g(x)|  < +\infty \]
\end{Example}

\begin{Proposition}{}{}
	Пусть $\mathcal{A} \colon X \rightarrow Y$~--- линейный, $\qquad X, Y$~--- банаховы пространства, тогда
	
	\[  \|\mathcal{A}\|  = \sup_{ \|x\| _X = 1} \|\mathcal{A}x\| _Y = \sup_{x \in X, x \neq 0} \frac{ \|\mathcal{A}x\| }{ \|x\| } \]
\end{Proposition}

\begin{proof}
	
	\[  \|\mathcal{A}\|  \ge \sup_{ \|x\| _X = 1}  \|\mathcal{A}x\| _Y \qquad (\text{на самом деле равенство})\]
	
	\[ \forall x \neq 0 \st  \|\mathcal{A}x\| _Y =  \|\mathcal{A}(\frac{x}{\| x \|_X})\| _Y\cdot  \|x\| _X \]
	
	\[ x = 0 \Rightarrow  \|\mathcal{A}x\| _Y = 0 \]
\end{proof}

\begin{Proposition}{}{}
	
	$X, Y$~--- банаховы пространства.
	\[
	\mathcal{A} \colon X \rightarrow Y \text{~--- ограничен} \Leftrightarrow \mathcal{A} \text{~--- непрерывный как отображение}  \Leftrightarrow \text{непрерывный в нуле.}
	\]
	
\end{Proposition}

\begin{proof}
	
	\begin{enumerate}
		\item Непрерывность в нуле $\Rightarrow$ непрерывность
		
		Неравенство в нуле:
		\[ \forall  \eps  >  0  \exists  \delta  >  0 \qquad  \|x\| _X  <  \delta \Rightarrow  \|\mathcal{A}x\| _Y  <  \eps \]
		
		\item Непрерывность в точке $Y \in \operatorname{Dom} \mathcal{A}$
		
		\[ \forall  \eps  >  0  \exists  \delta  >  0 \qquad  \|y-x\| _X  <  \delta \Rightarrow  \|\mathcal{A}y - \mathcal{A}x\|  =  \|\mathcal{A}(y-x)\|   <  \eps \]
		
		\item Непрерывность $\Rightarrow$ ограниченность
		
		Покажем, что если $\mathcal{A}$~--- непрерывный, то $ \|\mathcal{A}\|   <  + \infty$
		
		Рассмотрим $\sup_{ \|x\| _X \le 1}  \|\mathcal{A}x\| _Y$
		
		Непрерывность в нуле: 
		
		\[ \forall  \eps  >  0  \exists  \delta  >  0 \st  \|x\| _X \le \delta \Rightarrow  \|\mathcal{A}x\| _Y \le \eps \]
		
		Положим $\eps := 1$. Если $\|x'\|_X \le 1$, то $\|\delta x'\|_X \le \delta$, поэтому
		
		\[ \forall x' \st \|x'\|_X \le 1 \quad \|\mathcal{A}x'\|_Y = \frac{1}{\delta}\|\mathcal{A}(\delta x')\|_Y \le \frac{1}{\delta}. \]
		
		\[ \sup_{ \|x'\|_X \le 1}  \|\mathcal{A}x'\| _Y \le \frac{1}{\delta} \Rightarrow  \|\mathcal{A}\|  \le \frac{1}{\delta} \]
		
		\item Ограниченность $\Rightarrow$ непрерывность в нуле
		
		\[  \|\mathcal{A} x\|_Y \le  \|\mathcal{A} \|\cdot\| x \|_X  \forall  x \in X \]
		\[ \|\mathcal{A}\| = \sup_{X \ni x \neq 0} \frac{\|\mathcal{A}x\|}{\|x\|} \]
	\end{enumerate}
\end{proof}

\begin{Theorem}
	Пусть $X$ --- нормированное пространство, а $Y$ --- банахово пространство.
	Тогда пространство $\mathcal B(X,Y)$ ограниченных линейных операторов
	$A:X\to Y$ является банаховым пространством относительно операторной нормы
	\[
	\|A\|=\sup_{x\ne 0}\frac{\|Ax\|_Y}{\|x\|_X}.
	\]
\end{Theorem}

\begin{proof}
	Докажем полноту $\mathcal B(X,Y)$.
	
	Пусть $(A_n)$ --- фундаментальная последовательность в $\mathcal B(X,Y)$.
	Тогда для любого $\varepsilon>0$ существует $N\in\mathbb N$ такое, что при
	$n,m\ge N$
	\[
	\|A_n-A_m\|<\varepsilon.
	\]
	Зафиксируем $x\in X$. Тогда
	\[
	\|A_nx-A_mx\|_Y
	=
	\|(A_n-A_m)x\|_Y
	\le
	\|A_n-A_m\|\|x\|_X
	<
	\varepsilon\|x\|_X.
	\]
	Значит, $(A_nx)$ --- фундаментальная последовательность в $Y$.
	Так как $Y$ банахово, она сходится. Определим оператор $A:X\to Y$:
	\[
	Ax:=\lim_{n\to\infty}A_nx.
	\]
	
	Проверим линейность. Для $x,y\in X$ и $\lambda\in\mathbb K$ имеем
	\[
	A(\lambda x+y)
	=
	\lim_{n\to\infty}A_n(\lambda x+y)
	=
	\lim_{n\to\infty}(\lambda A_nx+A_ny)
	=
	\lambda Ax+Ay.
	\]
	Следовательно, $A$ линеен.
	
	Покажем, что $A_n\to A$ по операторной норме. При $n,m\ge N$
	\[
	\|(A_n-A_m)x\|_Y<\varepsilon\|x\|_X.
	\]
	Переходя к пределу при $m\to\infty$, получаем
	\[
	\|(A_n-A)x\|_Y\le \varepsilon\|x\|_X.
	\]
	Поэтому для $x\ne 0$
	\[
	\frac{\|(A_n-A)x\|_Y}{\|x\|_X}\le \varepsilon.
	\]
	Беря супремум по всем $x\ne 0$, получаем
	\[
	\|A_n-A\|\le \varepsilon.
	\]
	Значит,
	\[
	A_n\to A
	\]
	в операторной норме.
	
	Осталось показать, что $A\in\mathcal B(X,Y)$. Возьмём $n\ge N$. Тогда
	\[
	A=A_n-(A_n-A).
	\]
	Оператор $A_n$ ограничен, а оператор $A_n-A$ ограничен, так как
	$\|A_n-A\|\le\varepsilon$. Следовательно, $A$ тоже ограничен.
	
	Итак, $A$ --- ограниченный линейный оператор, причём $A_n\to A$ в
	$\mathcal B(X,Y)$. Значит, всякая фундаментальная последовательность в
	$\mathcal B(X,Y)$ сходится, то есть $\mathcal B(X,Y)$ является банаховым
	пространством.
\end{proof}